%!TEX root = ../main.tex
% GLOSARIO: términos que necesitan definición dentro del documento.
%
% Definir:   \newglossaryentry{clave}{name={Término}, description={Definición}}
% Usar:      \gls{clave}      -> Término (con enlace al glosario)
% main.tex usa \glsaddall, por lo que se imprimen todas las entradas.

\newglossaryentry{accesibilidad}{
    name={accesibilidad (reachability)},
    description={Propiedad de un grafo dirigido que indica si existe una
    secuencia de aristas que conecte un nodo de origen con un nodo de
    destino}
}

\newglossaryentry{taint}{
    name={análisis de contaminación (taint analysis)},
    description={Método que verifica si datos provenientes de fuentes no
    confiables pueden alcanzar operaciones sensibles sin una sanitización
    intermedia}
}

\newglossaryentry{deuda}{
    name={deuda técnica de seguridad},
    description={Costo diferido y riesgo acumulado en un sistema por
    decisiones de diseño o deficiencias de seguridad no resueltas en el
    código fuente}
}

\newglossaryentry{exhaustividad}{
    name={exhaustividad (recall)},
    description={Proporción de vulnerabilidades reales detectadas sobre el
    total de vulnerabilidades existentes}
}

\newglossaryentry{exposicion}{
    name={exposición al riesgo},
    description={Producto de la probabilidad de ocurrencia de un riesgo por
    su impacto. Se usa para ordenar la matriz de riesgos y decidir cuáles
    se tratan primero}
}

\newglossaryentry{fn}{
    name={falso negativo},
    description={Vulnerabilidad real presente en el código fuente que el
    analizador no detecta}
}

\newglossaryentry{fp}{
    name={falso positivo},
    description={Alerta emitida por el analizador sobre un fragmento de
    código seguro o correctamente sanitizado}
}

\newglossaryentry{dfgterm}{
    name={grafo de flujo de datos},
    description={Grafo dirigido cuyos vértices representan variables,
    operaciones e identificadores, y cuyas aristas modelan la propagación
    de los valores}
}

\newglossaryentry{precision}{
    name={precisión (precision)},
    description={Proporción de vulnerabilidades reales sobre el total de
    alertas emitidas por el analizador}
}

\newglossaryentry{prototipo}{
    name={prototipo},
    description={Versión preliminar de un sistema construida para validar
    una idea, un flujo o una interfaz. Puede simular partes del
    comportamiento y no está pensada para operar en producción}
}

\newglossaryentry{sanitizer}{
    name={sanitizer},
    description={Operación que valida, filtra o codifica un dato y le quita
    la condición de contaminado}
}

\newglossaryentry{sink}{
    name={sink},
    description={Operación sensible en la que el uso de un dato no confiable
    produce una vulnerabilidad, como una consulta a una base de datos o la
    ejecución de un comando}
}

\newglossaryentry{source}{
    name={source},
    description={Punto por el que ingresa al programa información externa
    no controlada, como los parámetros de una petición HTTP}
}

\newglossaryentry{vulnarq}{
    name={vulnerabilidad de arquitectura},
    description={Defecto de seguridad que no se manifiesta en una sentencia
    aislada sino en la forma en que los componentes se vinculan y en el
    recorrido de la información entre ellos}
}

\newglossaryentry{g:mvp}{
    name={producto mínimo viable},
    description={Versión funcional de un producto con el conjunto mínimo
    de características que entrega valor real a los usuarios y permite
    aprender de su uso. A diferencia de un prototipo, no simula: funciona}
}
